\chapter*{ТЕОРЕТИЧЕСКАЯ ЧАСТЬ}
\addcontentsline{toc}{chapter}{ТЕОРЕТИЧЕСКАЯ ЧАСТЬ}

\section*{Вопросы из лабораторной 5}

\subsection*{1. Виды симметричного шифрования (поточные и блочные). Схема для одного из видов}

Симметричное шифрование — это метод криптографии, в котором для шифрования и расшифровки используется один и тот же секретный ключ. Существуют два основных вида симметричного шифрования:

\subsubsection*{1.1. Поточное (потоковое) шифрование}

\textbf{Поточное шифрование} обрабатывает данные побитово или побайтово. Генератор создаёт поток ключевой последовательности, которая складывается с исходным текстом операцией XOR:
$$C_i = M_i \oplus K_i$$

где $C_i$ — шифротекст, $M_i$ — открытый текст, $K_i$ — ключевой поток.

\textbf{Примеры поточных шифров:} RC4, ChaCha20, Salsa20, A5/1.

\textbf{Преимущества:}
\begin{itemize}
    \item Высокая скорость работы
    \item Низкие требования к памяти
    \item Подходит для потоковых данных (аудио, видео)
\end{itemize}

\textbf{Недостатки:}
\begin{itemize}
    \item Критична синхронизация потока
    \item Повторное использование ключа крайне опасно
\end{itemize}

\textbf{Схема поточного шифрования:}

\begin{verbatim}
Открытый текст → [XOR] ← Ключевой поток
                    ↓
              Шифротекст
\end{verbatim}

\subsubsection*{1.2. Блочное шифрование}

\textbf{Блочное шифрование} обрабатывает данные блоками фиксированного размера (обычно 64 или 128 бит). Исходный текст разбивается на блоки, каждый блок проходит через серию раундов с подстановками и перестановками.

\textbf{Примеры блочных шифров:} DES (64-битные блоки), AES (128-битные блоки), 3DES, ГОСТ 28147-89.

\textbf{Преимущества:}
\begin{itemize}
    \item Высокая криптостойкость
    \item Возможность параллельной обработки
    \item Хорошо стандартизированы
\end{itemize}

\textbf{Недостатки:}
\begin{itemize}
    \item Требуется дополнение данных до размера блока (padding)
    \item Необходимость выбора режима работы (ECB, CBC, CTR и др.)
\end{itemize}

\textbf{Схема блочного шифрования (режим ECB):}

\begin{verbatim}
Блок 1 → [Шифрование] → Шифрованный блок 1
           с ключом

Блок 2 → [Шифрование] → Шифрованный блок 2
           с ключом

Блок N → [Шифрование] → Шифрованный блок N
           с ключом
\end{verbatim}

\textit{В данной работе реализован блочный шифр AES с размером блока 128 бит (16 байт).}

\subsection*{2. Описание алгоритма шифрования AES}

\textbf{AES (Advanced Encryption Standard)} — симметричный алгоритм блочного шифрования, принятый в качестве стандарта правительством США в 2001 году. AES использует 128-битные блоки данных и поддерживает ключи длиной 128, 192 или 256 бит.

\subsubsection*{Структура алгоритма AES}

Алгоритм AES состоит из следующих этапов:

\begin{enumerate}
    \item \textbf{Key Expansion (Расширение ключа):} Из исходного ключа генерируются раундовые ключи для каждого раунда шифрования. Для AES-256 генерируется 15 раундовых ключей (60 32-битных слов).
    
    \item \textbf{Initial Round (Начальный раунд):} Применяется операция AddRoundKey — XOR блока данных с первым раундовым ключом.
    
    \item \textbf{Main Rounds (Основные раунды):} Для AES-256 выполняется 13 основных раундов (раунды 1-13). Каждый раунд включает четыре операции:
    \begin{itemize}
        \item \textbf{SubBytes}: Нелинейная подстановка байтов через S-блок
        \item \textbf{ShiftRows}: Циклический сдвиг строк матрицы состояния
        \item \textbf{MixColumns}: Перемешивание столбцов матрицы состояния
        \item \textbf{AddRoundKey}: XOR с раундовым ключом
    \end{itemize}
    
    \item \textbf{Final Round (Финальный раунд):} 14-й раунд выполняется без операции MixColumns:
    \begin{itemize}
        \item SubBytes
        \item ShiftRows
        \item AddRoundKey
    \end{itemize}
\end{enumerate}

\subsubsection*{Расшифровка AES}

Расшифровка выполняется тем же алгоритмом с использованием обратных операций (InvSubBytes, InvShiftRows, InvMixColumns) и раундовых ключей в обратном порядке.

\subsubsection*{Генерация раундовых ключей (Key Expansion)}

Для AES-256:
\begin{itemize}
    \item Исходный 256-битный ключ делится на 8 32-битных слов
    \item Всего генерируется 60 32-битных слов (4 слова $\times$ 15 раундов)
    \item Используются операции RotWord (циклический сдвиг), SubWord (подстановка через S-блок) и XOR с константами Rcon
\end{itemize}

\subsection*{3. Определения алгоритмов перестановки и подстановки. Примеры каждого из этих видов алгоритмов. Пример алгоритма, использующего оба подхода}

\subsubsection*{3.1. Перестановка (Permutation)}

\textbf{Перестановка} — это криптографическое преобразование, которое переставляет биты или байты входной последовательности в определённом порядке согласно таблице перестановок.

Для входа $X = x_1x_2...x_n$ и таблицы перестановок $P = [p_1, p_2, ..., p_m]$ выход определяется как:
$$y_i = x_{p_i}, \quad i = 1, 2, ..., m$$

\textbf{Свойства перестановки:}
\begin{itemize}
    \item Линейность
    \item Обратимость
    \item Обеспечивает диффузию (распространение влияния каждого бита)
\end{itemize}

\textbf{Примеры перестановок в криптографии:}
\begin{itemize}
    \item Начальная перестановка (IP) в DES
    \item Расширение E в DES ($32 \rightarrow 48$ бит)
    \item Перестановка P после S-блоков в DES
    \item ShiftRows в AES
\end{itemize}

\subsubsection*{3.2. Подстановка (Substitution)}

\textbf{Подстановка} — это нелинейное преобразование, заменяющее входную последовательность на другую согласно таблице замен (S-box):
$$y = S(x)$$

\textbf{Свойства подстановки:}
\begin{itemize}
    \item Нелинейность
    \item Обеспечивает конфузию (скрывает связь между ключом и шифротекстом)
    \item Устойчивость к линейному и дифференциальному криптоанализу
\end{itemize}

\textbf{Примеры подстановок в криптографии:}
\begin{itemize}
    \item S-блоки в DES (8 таблиц, $6$ бит $\rightarrow$ $4$ бита)
    \item S-box в AES (таблица $16 \times 16$, $8$ бит $\rightarrow$ $8$ бит)
    \item Шифр Цезаря (простейшая подстановка)
\end{itemize}

\textbf{SubBytes в AES:}

В AES операция SubBytes заменяет каждый байт матрицы состояния на соответствующее значение из S-блока. S-блок построен на основе обратного элемента в поле Галуа $GF(2^8)$ с последующим аффинным преобразованием.

\subsubsection*{3.3. Алгоритм, использующий оба подхода}

\textbf{AES} является классическим примером SP-сети (Substitution-Permutation Network) — алгоритма, комбинирующего подстановки и перестановки:

\textbf{Подстановки в AES:}
\begin{itemize}
    \item SubBytes — нелинейная подстановка через S-блок, обеспечивает конфузию
\end{itemize}

\textbf{Перестановки в AES:}
\begin{itemize}
    \item ShiftRows — циклический сдвиг строк матрицы состояния
    \item MixColumns — линейное преобразование столбцов (можно рассматривать как обобщённую перестановку)
    \item Key Expansion использует перестановки (RotWord)
\end{itemize}

\textbf{Принцип работы:}

Чередование операций подстановки и перестановки в 14 раундах (для AES-256) создаёт высокую криптографическую стойкость:
\begin{itemize}
    \item \textbf{Перестановки} обеспечивают \textbf{диффузию} — распространение влияния каждого бита входа на весь блок выхода
    \item \textbf{Подстановки} обеспечивают \textbf{конфузию} — усложнение статистической связи между ключом и шифротекстом
\end{itemize}

Благодаря многократному применению этих операций, изменение даже одного бита входных данных приводит к изменению в среднем половины бит выходных данных (лавинный эффект).

\section*{Вопросы из лабораторной 6}

\subsection*{1. Классификация автоматизированных систем с точки зрения требований безопасности (3 класса защищённости)}

Автоматизированные системы классифицируются по требованиям безопасности на основе важности обрабатываемой информации и потенциальных угроз:

\textbf{Класс 3 (низкий уровень защищённости):}
\begin{itemize}
    \item Обрабатывается информация ограниченного распространения
    \item Базовые требования: аутентификация пользователей, журналирование событий, антивирусная защита
    \item Пример: корпоративные информационные системы без критичных данных
\end{itemize}

\textbf{Класс 2 (средний уровень защищённости):}
\begin{itemize}
    \item Обрабатывается конфиденциальная информация
    \item Дополнительные требования: контроль доступа, шифрование данных, резервное копирование, защита от несанкционированного доступа
    \item Пример: банковские системы, медицинские информационные системы
\end{itemize}

\textbf{Класс 1 (высокий уровень защищённости):}
\begin{itemize}
    \item Обрабатывается информация, составляющая государственную тайну
    \item Строгие требования: многоуровневая аутентификация, криптографическая защита всех каналов, физическая защита оборудования, сертифицированные средства защиты
    \item Пример: государственные информационные системы специального назначения
\end{itemize}

\subsection*{2. Виды симметричного шифрования (поточные и блочные)}

Данный вопрос был подробно рассмотрен в Блоке 1, пункт 1.

\subsection*{3. Особенности алгоритма шифрования архивных файлов}

При шифровании архивных файлов необходимо учитывать следующие особенности:

\subsubsection*{3.1. Структура архивных файлов}

Архивные файлы (ZIP, RAR, 7z) имеют внутреннюю структуру:
\begin{itemize}
    \item Заголовки с метаданными
    \item Сжатое содержимое файлов
    \item Каталог архива
\end{itemize}

\subsubsection*{3.2. Подходы к шифрованию архивов}

\textbf{Шифрование на уровне архиватора:}
\begin{itemize}
    \item Многие архиваторы поддерживают встроенное шифрование (например, ZIP с AES)
    \item Шифруется только содержимое файлов, метаданные остаются открытыми
    \item Преимущество: возможность просмотра списка файлов без расшифровки
\end{itemize}

\textbf{Шифрование архива целиком (внешнее шифрование):}
\begin{itemize}
    \item Архив рассматривается как обычный бинарный файл
    \item Шифруется весь файл, включая заголовки и метаданные
    \item Преимущество: максимальная конфиденциальность (невозможно определить даже структуру архива)
    \item Недостаток: для доступа к любому файлу требуется полная расшифровка
\end{itemize}

\subsubsection*{3.3. Особенности реализации}

В данной работе используется подход внешнего шифрования:
\begin{itemize}
    \item Архивный файл обрабатывается как последовательность байт
    \item Применяется блочное шифрование AES-256 в режиме ECB
    \item Zero padding для последнего блока
    \item Сохранение исходного размера файла для корректной расшифровки
\end{itemize}

Преимущества такого подхода:
\begin{itemize}
    \item Универсальность (подходит для любых файлов)
    \item Простота реализации
    \item Полная конфиденциальность (невозможно определить тип файла)
\end{itemize}

\section*{Вопросы из лабораторной 7}

\subsection*{1. Перечислите подсистемы защиты информации в автоматизированных системах}

Основные подсистемы защиты информации в автоматизированных системах:

\subsubsection*{1.1. Подсистема идентификации и аутентификации}
\begin{itemize}
    \item Проверка подлинности пользователей и процессов
    \item Методы: пароли, биометрия, электронные ключи, двухфакторная аутентификация
\end{itemize}

\subsubsection*{1.2. Подсистема управления доступом}
\begin{itemize}
    \item Контроль прав доступа к ресурсам системы
    \item Модели: дискреционная (DAC), мандатная (MAC), ролевая (RBAC)
\end{itemize}

\subsubsection*{1.3. Подсистема шифрования данных}
\begin{itemize}
    \item Защита информации при хранении и передаче
    \item Симметричное и асимметричное шифрование, криптографические протоколы
\end{itemize}

\subsubsection*{1.4. Подсистема регистрации и аудита}
\begin{itemize}
    \item Журналирование событий безопасности
    \item Анализ действий пользователей и системных процессов
\end{itemize}

\subsubsection*{1.5. Подсистема антивирусной защиты}
\begin{itemize}
    \item Обнаружение и нейтрализация вредоносного ПО
    \item Регулярное обновление антивирусных баз
\end{itemize}

\subsubsection*{1.6. Подсистема резервного копирования}
\begin{itemize}
    \item Создание копий критичных данных
    \item Обеспечение восстановления после сбоев или атак
\end{itemize}

\subsubsection*{1.7. Подсистема межсетевого экранирования}
\begin{itemize}
    \item Фильтрация сетевого трафика
    \item Защита от внешних и внутренних угроз
\end{itemize}

\subsection*{2. Описание алгоритма шифрования AES}

Данный вопрос был подробно рассмотрен в Блоке 1, пункт 2.

\subsection*{3. Особенности алгоритма симметричного шифрования архивных файлов}

Данный вопрос был подробно рассмотрен в Блоке 2, пункт 3.